\subsection{Generování mezikódu}
\label{sec:ir}

Mezikód (IR) je generován v~modulu \texttt{src/ir\_lowering.rs} transformací
AST na strukturu \texttt{IRProgram}. Implementace používá pomocnou strukturu
\texttt{IRBuilder} (\texttt{src/ir.rs}), která umožňuje postupně emitovat
instrukce do právě budované funkce.

Lowering je realizován ve třech krocích, které odpovídají přirozené struktuře
\Forth{}-like programu:
\begin{enumerate}
	\item nejprve se sesbírají všechny definice slov do mapy \texttt{word\_definitions},
	\item následně se pro každé slovo vytvoří samostatná IR funkce,
	\item nakonec je přeložena hlavní část programu (kód mimo definice) do funkce \texttt{main}.
\end{enumerate}

Číselné literály jsou překládány na instrukci \texttt{Push(Constant)}.
Řetězcové literály jsou v~současné implementaci reprezentovány jako sekvence
znaků uložených na zásobník (ASCII kódy) následovaná délkou řetězce; to
odpovídá očekávanému stack efektu slova \texttt{TYPE}.

\paragraph{Proměnné.}
Při zpracování uzlu \texttt{VariableDeclaration} je proměnné přidělena celočíselná
adresa z~rostoucího čítače. Při použití proměnné jako slova (tj. token \texttt{Word}
odpovídá jménu proměnné) lowering vloží na zásobník její adresu.
Operace \texttt{@} a~\texttt{!} jsou mapovány na instrukce \texttt{Load} a~\texttt{Store}.

\paragraph{Řídicí struktury.}
Základní struktury \texttt{IF/ELSE/THEN} a~\texttt{DO/LOOP} jsou překládány na
instrukce se skoky a~návěstími (\texttt{Jump}, \texttt{JumpIfNot}, \texttt{Label})
pomocí generování unikátních návěstí v~\texttt{IRBuilder}. Pro správné párování
konstrukcí lowering udržuje pomocné zásobníky (\texttt{loop\_stack} a~\texttt{conditional\_stack}).

\paragraph{Výpočet stack effect.}
Po vygenerování IR je nad funkcemi spuštěna analýza stack efektu
(\texttt{StackEffectAnalyzer}), která z~anotací instrukcí (procedurální makro
\texttt{\#[derive(StackEffect)]} v~\texttt{roth-derive}) vypočítá čistou změnu
hloubky zásobníku pro každou funkci. Tato informace je později využita
optimalizátorem i~pro ladicí výpisy.
